\section{Strict positivity of original geometric ample numerical classes}
\label{sec:native-numerical-ample-classes}

\begin{theorem}[Geometric ampleness in the true original numerical quotient {\cite[\S1.2]{BJ22}}]
\label{thm:native-numerical-ample-classes}
\lean{EnvelopeProperty.NativeRealNumericalClasses.curveDegree_ampleLineClass_pos, EnvelopeProperty.NativeRealNumericalClasses.ampleLineClass_mem_nefCone, EnvelopeProperty.NativeRealNumericalClasses.ampleLineClass_ne_zero_of_curve, EnvelopeProperty.NativeRealNumericalClasses.curveDegree_positive_ample_sum, EnvelopeProperty.NativeRealNumericalClasses.positive_ample_sum_mem_nefCone}
\leanok
\uses{thm:native-numerical-nef-cone, thm:fc-original-geometric-ample-curve-degree-positive}
Let $X$ be projective over its original field. The genuine
numerical class of every actual geometric ample line has
strictly positive degree on every original closed integral
curve, and therefore belongs to the genuine closed-curve nef
cone. If an original closed integral curve exists, that class
is nonzero. Every finite nonempty positive real combination
of such original ample line classes has the same strict
positivity and belongs to the same nef cone.
\end{theorem}

\begin{proof}
The descended numerical degree of the original line class is
its true full Euler degree after pullback to the original
curve. Geometric ample curve positivity proves this integer
strictly positive; its image in the reals is strictly positive
as well. Nonnegative degree on every curve is exactly membership
in the original numerical nef cone. If the numerical class
were zero, its degree on a given original curve would be zero,
contradicting strict positivity.

For a positive real combination, linearity of each original
curve functional evaluates the class as the finite sum of
the coefficient times the corresponding positive original
curve degree. Every summand is positive, and the indexing
set is nonempty, so the sum is strictly positive. Applying
this to each original closed integral curve gives nef
membership. No finite-dimensionality or identification of
an ample topology is used in these statements.
\end{proof}
